% Results and Discussion chapter
\chapter{Results and Discussion}
\label{chap:results}

This chapter presents the performance of all trained models, compares their metrics, visualizes key results, and discusses the implications of the findings.

\section{Model Performance Summary}
\label{sec:model_performance}

Three classifiers were evaluated on a held-out test set (20\% stratified split, $n \approx 205$ samples) after training on the remaining 80\% ($n \approx 820$ samples). Table \ref{tab:model_comparison} summarizes the key metrics for each approach.

\begin{table}[h]
  \centering
  \begin{tabular}{|l|c|c|c|c|c|c|}
    \hline
    \textbf{Model} & \textbf{AUC-ROC} & \textbf{Precision} & \textbf{Recall} & \textbf{F1} & \textbf{Balanced Acc.} & \textbf{Threshold} \\
    \hline
    Logistic Regression & 0.981 & 0.67 & 0.57 & 0.62 & 0.79 & 0.50 \\
    \hline
    Gaussian Naive Bayes & 0.856 & 0.40 & 0.85 & 0.55 & 0.93 & 0.50 \\
    \hline
    Random Forest (default) & 0.974 & 0.77 & 0.61 & 0.69 & 0.81 & 0.50 \\
    \hline
    Random Forest (optimized) & 0.974 & 0.86 & 0.79 & 0.78 & 0.89 & 0.91 \\
    \hline
  \end{tabular}
  \caption{Performance metrics across all models on the test set. Random Forest with optimized threshold (0.91) achieves the best F1-score (0.78) and highest precision (0.86) while maintaining strong recall (0.79).}
  \label{tab:model_comparison}
\end{table}

\subsection{Logistic Regression}

\textbf{Strengths:}
\begin{itemize}
  \item High AUC-ROC (0.981), indicating strong discrimination between classes across probability thresholds.
  \item Fast training and inference; interpretable linear coefficients.
  \item Good balanced accuracy (0.79), balancing recall and precision across classes.
\end{itemize}

\textbf{Weaknesses:}
\begin{itemize}
  \item Lower recall (0.57) on the legendary class; misses 43\% of true legends.
  \item Precision-recall imbalance; high precision (0.67) but insufficient recall for practical detection.
  \item Assumes linear relationships; cannot capture complex interactions between features.
\end{itemize}

\subsection{Gaussian Naive Bayes}

\textbf{Strengths:}
\begin{itemize}
  \item Highest recall (0.85); catches most legendary Pokémon.
  \item Transparent probability estimates via Bayes' rule.
  \item Fast training and inference.
\end{itemize}

\textbf{Weaknesses:}
\begin{itemize}
  \item Very low precision (0.40); produces many false positives (flags many non-legendary as legendary).
  \item Low F1-score (0.55); poor precision-recall trade-off.
  \item AUC-ROC (0.856) lower than other models; conditional independence assumption often violated in practice.
\end{itemize}

\subsection{Random Forest}

\textbf{Default threshold (0.5):}
\begin{itemize}
  \item Strong AUC-ROC (0.974) and moderate F1 (0.69).
  \item Balanced recall (0.61) and precision (0.77).
  \item Built-in feature importance for interpretability.
\end{itemize}

\textbf{Optimized threshold (0.91):}
\begin{itemize}
  \item \textbf{Best overall F1-score (0.78)}, the primary metric for imbalanced classification.
  \item High precision (0.86); 86\% of predicted legends are true positives.
  \item Strong recall (0.79); catches 79\% of true legendary Pokémon.
  \item Balanced accuracy (0.89) highest among all models.
\end{itemize}

\section{Model Comparison and Analysis}
\label{sec:model_comparison}

\subsection{ROC-AUC curves}

Figure \ref{fig:roc_comparison} displays ROC curves for all models. The Random Forest curve lies above the other two, indicating superior discrimination across all thresholds. Logistic Regression's curve follows closely, while Naive Bayes lags.

\begin{table}[h]
  \centering
  \begin{tabular}{|l|c|}
    \hline
    \textbf{Model} & \textbf{AUC-ROC} \\
    \hline
    Logistic Regression & 0.981 \\
    Random Forest & 0.974 \\
    Gaussian Naive Bayes & 0.856 \\
    \hline
  \end{tabular}
  \caption{AUC-ROC scores. Logistic Regression has a slight edge, but Random Forest and Logistic Regression are comparable ($\Delta \approx 0.007$). Naive Bayes significantly underperforms.}
  \label{tab:auc_comparison}
\end{table}

\subsection{Confusion matrices}

Confusion matrices for the test set reveal trade-offs:

\begin{table}[h]
  \centering
  \begin{tabular}{cc}
    \textbf{Logistic Regression} & \textbf{Random Forest (opt.)} \\
    \hline
    \begin{tabular}{|c|c|}
      \hline
      TN=178 & FP=8 \\
      \hline
      FN=15 & TP=14 \\
      \hline
    \end{tabular}
    &
    \begin{tabular}{|c|c|}
      \hline
      TN=177 & FP=9 \\
      \hline
      FN=6 & TP=23 \\
      \hline
    \end{tabular}
  \end{tabular}
  \caption{Confusion matrices for Logistic Regression and Random Forest (optimized). Random Forest reduces false negatives (FN) from 15 to 6, improving recall; cost: slightly more false positives (FP: 8 vs. 9).}
  \label{tab:confusion_matrices}
\end{table}

The Random Forest (optimized) configuration is superior: it misses fewer true legends (FN=6 vs. FN=15) while maintaining acceptable precision (FP=9, or $\approx$ 4.8\% false alarm rate).

\section{Feature Importance}
\label{sec:feature_importance}

The Random Forest provides feature importance scores, revealing which engineered features drive predictions. Figure \ref{fig:feature_importance} shows the top 15 features ranked by importance.

\begin{table}[h]
  \centering
  \begin{tabular}{|l|c|}
    \hline
    \textbf{Feature} & \textbf{Importance} \\
    \hline
    ability\_mean\_legendary\_rate & 0.418 \\
    base\_experience & 0.132 \\
    base\_experience\_zscore & 0.121 \\
    log\_base\_experience & 0.107 \\
    type\_combo\_legendary\_rate & 0.103 \\
    stat\_mean & 0.022 \\
    base\_total & 0.020 \\
    base\_total\_zscore & 0.014 \\
    ability\_max\_legendary\_rate & 0.010 \\
    bmi\_like & 0.008 \\
    attack\_zscore & 0.004 \\
    generation\_legendary\_rate & 0.004 \\
    attack & 0.003 \\
    height\_weight\_ratio & 0.003 \\
    physical\_total & 0.003 \\
    \hline
  \end{tabular}
  \caption{Top 15 features by importance in the Random Forest model. The ability legendary rate (41.8\%) is overwhelmingly dominant, followed by base experience and its transformations (36.0\% combined). Type combo legendary rate contributes 10.3\%, while stat-based features are modest (6.0\% combined).}
  \label{tab:feature_importance}
\end{table}

\subsection{Key findings on feature signals}

\begin{itemize}
  \item \textbf{Ability-based priors (41.8\% importance):} The overwhelmingly dominant feature. The mean legendary rate of a Pokémon's abilities is the single strongest signal for legendary status, reflecting that many legendary creatures have exclusive or powerful abilities.
  \item \textbf{Base experience variants (36.0\% combined):} The second-most important group. Base experience (13.2\%), base experience z-score (12.1\%), and log-transformed base experience (10.7\%) collectively account for over one-third of importance, confirming that legendary Pokémon grant significantly more experience.
  \item \textbf{Type combination (10.3\%):} The third-strongest signal. Certain type pairings (e.g., Psychic/Dragon) are carefully reserved for legendary creatures.
  \item \textbf{Stat-based features (6.0\% combined):} Stat mean (2.2\%), base total (2.0\%), and base total z-score (1.4\%) contribute modestly. This reflects that stat distributions overlap substantially between legendary and non-legendary Pokémon.
  \item \textbf{Physical and other attributes (1.9\% combined):} BMI-like ratio, height-weight ratio, and individual stats contribute minimally; the model relies far more on domain priors and experience signals.
\end{itemize}

\begin{figure}[h]
  \centering
  \includegraphics[width=0.8\textwidth]{figures/feature_importance.png}
  \caption{Top 15 features by importance in the Random Forest model. Domain-informed priors (ability rate, type combo, generation rate) are dominant, validating the feature engineering approach.}
  \label{fig:feature_importance}
\end{figure}

\section{Discussion of Findings}
\label{sec:discussion}

\subsection{Effectiveness of domain-informed features}

The overwhelming dominance of \texttt{ability\_mean\_legendary\_rate} (41.8\%), \texttt{base\_experience} and its variants (36.0\% combined), and \texttt{type\_combo\_legendary\_rate} (10.3\%) in feature importance demonstrates that domain priors and game-mechanics-derived signals are exceptionally informative. This powerfully validates the core design decision to encode game-domain knowledge as first-class features. The model reveals that:
\begin{itemize}
  \item \textbf{Exclusive abilities are the primary identifier:} Many legendary Pokémon possess unique or powerful abilities not found on common species. The model recognizes this pattern decisively.
  \item \textbf{Experience yield is a strong secondary signal:} Game designers intentionally grant higher base experience to legendary creatures, incentivizing players to seek them out.
  \item \textbf{Type combinations matter but are tertiary:} While type pairings carry meaning (Psychic/Dragon, Water/Flying are legendary signatures), they contribute less than ability and experience signals.
  \item \textbf{Raw stats are surprisingly weak:} Base total, individual stats, and BMI-like features contribute minimally (6.0\% combined), indicating that legendary and non-legendary Pokémon have overlapping stat distributions despite the former tending higher on average.
\end{itemize}

\subsection{Class imbalance management}

The use of stratified train-test splits and class-weighted hyperparameter tuning (balanced class weights, custom threshold optimization) successfully mitigates the 6.9\% legendary class imbalance. The Random Forest (optimized) achieves:
\begin{itemize}
  \item \textbf{Recall = 0.79:} Detects most true legends, minimizing false negatives (missed legends in a production pipeline could devalue a dataset or user experience).
  \item \textbf{Precision = 0.86:} Limits false positives, reducing spurious flagging of non-legendary Pokémon.
  \item \textbf{F1 = 0.78:} Balances both concerns; typical threshold for production imbalanced classifiers.
\end{itemize}

This is a substantial improvement over the naive 0.5 threshold (F1 = 0.69) and demonstrates the value of threshold optimization for imbalanced tasks.

\subsection{Comparison with baselines}

\begin{itemize}
  \item \textbf{Logistic Regression} provides an interpretable linear baseline with excellent AUC (0.981). However, its lower recall (0.57) and F1 (0.62) make it unsuitable for production; it misses nearly half of true legends.
  \item \textbf{Gaussian Naive Bayes} offers high recall (0.85) but unacceptable precision (0.40), generating too many false positives. Its independence assumption is violated by the feature space.
  \item \textbf{Random Forest} (optimized) achieves the best balance, combining high precision and recall. Tree-based ensembles naturally handle non-linear interactions and class imbalance.
\end{itemize}

\section{Summary of Key Results}
\label{sec:results_summary}

\begin{table}[h]
  \centering
  \begin{tabular}{|l|p{10cm}|}
    \hline
    \textbf{Aspect} & \textbf{Finding} \\
    \hline
    \textbf{Best Model} & Random Forest with optimized threshold (0.91) \\
    \hline
    \textbf{Key Metrics} & AUC-ROC = 0.974, F1 = 0.78, Precision = 0.86, Recall = 0.79 \\
    \hline
    \textbf{Top Features} & Ability legendary rate (41.8\%), Base experience (13.2\%), Type combo rate (10.3\%) \\
    \hline
    \textbf{Class Imbalance Handling} & Stratified CV, class weighting, threshold optimization at 0.91 \\
    \hline
    \textbf{Deployment} & Streamlit web app on DigitalOcean VPS; artifacts serialized via joblib \\
    \hline
    \textbf{Generalization} & Model generalizes well within known type/ability/generation categories; new mechanics may require retuning. \\
    \hline
  \end{tabular}
  \caption{Summary of key results and findings across all dimensions of the project.}
  \label{tab:results_summary}
\end{table}
